%% CV for Norlys Energy Trading - Software Developer (Dansk)
%% Compiled with: lualatex

\documentclass[11pt,a4paper,sans]{moderncv}
\moderncvstyle{banking}
\moderncvcolor{blue}

\renewcommand*{\firstnamestyle}[1]{{\fontsize{34}{36}\bfseries\upshape\color{color1}#1}}
\renewcommand*{\lastnamestyle}[1]{{\fontsize{34}{36}\bfseries\upshape\color{color1}#1}}
\renewcommand*{\sectionstyle}[1]{{\sectionfont\color{color1}#1}}

\usepackage[utf8]{inputenc}
\usepackage{hyperref}
\hypersetup{
    colorlinks=true,
    linkcolor=blue,
    filecolor=magenta,
    urlcolor=blue,
    pdftitle={Jacob El-Omar - CV},
    pdfpagemode=FullScreen,
}
\usepackage[scale=0.80]{geometry}
\usepackage{import}

\name{Jacob}{El-Omar}
\address{Aalborg, Danmark}{}{}
\phone[mobile]{+45 53 35 27 82}
\email{elomarjc@gmail.com}
\extrainfo{\href{https://linkedin.com/in/jacob-el-omar}{LinkedIn} \textbullet{} \href{https://github.com/elomarjc}{GitHub}}

\begin{document}

\makecvtitle

% ============================================================
%     PROFIL
% ============================================================

\vspace{6pt}
\small{Nyuddannet civilingeniør (cand.polyt.) i elektroniske systemer fra Aalborg Universitet med speciale i softwareudvikling, distribuerede systemer og signalbehandling. Erfaren med C\#, SQL-databaser og Flutter-appudvikling samt et stærkt teoretisk og praktisk fundament inden for design af hændelsesstyret logik. En samarbejdsorienteret udvikler trænet i problembaseret læring (PBL), klar til at udvikle og optimere pålidelige cloud-baserede softwaretjenester i jeres Commercial Solutions Development-team.}

% ============================================================
%     FAGLIGE KOMPETENCER
% ============================================================

\section{Faglige Kompetencer}
\vspace{1pt}
\begin{itemize}

\item \textbf{Softwareudvikling \& API'er}: C\#, .NET Core, Python, JavaScript, TypeScript, Node.js, REST API'er, JSON.

\item \textbf{Databaser \& Integration}: PostgreSQL, MySQL, SQLite, Supabase, SQL-forespørgsler og relationel schema-optimering.

\item \textbf{Distribuerede \& Hændelsesstyrede systemer}: Trænet i stokastiske systemer, sensorkommunikation, hændelsessporing og integration via REST.

\item \textbf{Mobil- \& Webudvikling}: Flutter, Dart, WordPress-plugin-konfiguration og web-optimering.

\item \textbf{Ingeniørværktøjer}: git, GitHub Actions (CI/CD-workflows), Scrum-metodologi, Trello og teams-baseret projektstyring.

\end{itemize}

% ============================================================
%     ERHVERVSERFARING
% ============================================================

\section{Erhvervserfaring}
\vspace{3pt}
\begin{itemize}

\item{\cventry{Jan 2026--Nu}{IT- \& Webudvikler (Frivillig)}{Dawah Danmark}{Aalborg, Danmark}{}{\vspace{1pt}
\begin{itemize}
    \item Moderniserede WordPress-arkitekturen for at forbedre server-side ydeevne og indlæsningstider.
    \item Konfigurerede plugin-integrationer og optimerede databasedesign for bedre datakonsistens.
\end{itemize}}}

\vspace{3pt}

\item{\cventry{Sep 2024--Jun 2025}{AI/ML Engineer (Praktik)}{Ai Health Highway}{Remote}{}{\vspace{1pt}
\begin{itemize}
    \item Udviklede og implementerede signalbehandlingsalgoritmer i Python til støjreduktion i smarte stetoskoper.
    \item Udarbejdede og testede LMS/RLS adaptive filtre og FastICA-modeller i tæt samarbejde med internationale teams.
    \item Dokumenterede systemkrav og testprocedurer i overensstemmelse med standarder for medicinsk software.
\end{itemize}}}

\vspace{3pt}

\item{\cventry{Jun 2023--Aug 2024}{Lager- og Logistikmedarbejder}{Lyn Express}{Aalborg, Danmark}{}{\vspace{1pt}
\begin{itemize}
    \item Styrede sortering og registrering af forsendelser under tidspres i et travlt logistikmiljø.
    \item Automatiserede regnearksarbejdsgange for at optimere lagerstyring og teamkoordinering.
\end{itemize}}}

\end{itemize}

% ============================================================
%     UDDANNELSE
% ============================================================

\section{Uddannelse}
\vspace{1pt}
\begin{itemize}

\item{\cventry{Sep 2023--Jun 2025}{Civilingeniør, cand.polyt. (Elektroniske Systemer)}{Aalborg Universitet}{Aalborg, Danmark}{}{\vspace{1pt}
Speciale: ``Adaptive Noise Cancellation For Electronic Stethoscopes.'' Fokus på maskinlæring, signalbehandling, stokastiske systemer og distribueret software.
}}

\vspace{3pt}

\item{\cventry{Sep 2020--Aug 2023}{Bachelor i teknisk videnskab (Elektronik og IT)}{Aalborg University}{Aalborg, Danmark}{}{\vspace{1pt}
Specialisering i reguleringsteknik. Projekter omfattede bl.a. kollisionsforebyggende navigationssystemer til mobile robotter (MiR200) baseret på UWB og regressionsalgoritmer.
}}

\end{itemize}

% ============================================================
%     UDVALGTE PROJEKTER
% ============================================================

\section{Udvalgte Projekter}
\vspace{1pt}
\begin{itemize}

\item{\textbf{Grow a Fish (Fritidsprojekt)}: Udviklet med Flutter/Dart og en Supabase PostgreSQL-backend. Implementerede synkronisering i realtid og komplekse SQL-forespørgsler samt API-integrationer.}

\item{\textbf{Kollisionsforebyggelse til Mobile Robotter (Universitetsprojekt, 5. sem.)}: Udviklede et prædiktivt navigations- og kollisionsforebyggelsessystem til MiR200 mobile robotplatforme i realtid baseret på regressionsalgoritmer og UWB tracking-data.}

\item{\textbf{Embedded Faldsikringssystem (Universitetsprojekt, 4. sem.)}: Udviklede software til faldregistrering på en Cypress PSoC-mikrocontroller i C. Implementerede I2C-kommunikation, sensorfusion og overførsel af data.}

\end{itemize}

% ============================================================
%     SPROG
% ============================================================

\section{Sprog}
\vspace{1pt}
\begin{itemize}
\item Dansk (modersmål), Engelsk (flydende).
\end{itemize}

% ============================================================
%     REFERENCER
% ============================================================

\section{Referencer}
\vspace{1pt}
\begin{itemize}
\item{Afleveres efter aftale.}
\end{itemize}

\end{document}
